\chapter{Introduction}
\label{chap:introduction}
\minitoc

\section{Motivation}
\label{sec:introduction-motivation}

Minimally invasive surgery (MIS) has transformed surgical care by reducing tissue trauma, blood loss, post-operative pain, recovery time, and complication rates compared with open surgery \cite{jaffray2005minimally}. These benefits arise partly because the operative field is accessed through small incisions and viewed through an endoscope or laparoscope. However, the same visual mediation that makes MIS less invasive also makes the operative scene difficult to interpret computationally. The camera view is narrow, indirect, and continuously changing, while smoke, blood, specular reflections, motion blur, occlusion, tissue deformation, and variable illumination can obscure important visual evidence.

Computer-assisted intervention seeks to extract useful information from these complex visual streams. Potential applications include surgical workflow analysis, skill assessment, context-aware assistance, robotic support, and automatic documentation. All of these applications depend on a fundamental capability: a computational system must represent not only which objects are visible, but also what is happening between them.

Substantial progress has been made on individual surgical vision tasks. Models can recognise instruments, segment their visible regions, identify operative phases, and classify surgical activities. Nevertheless, surgical scene understanding requires more than a collection of independent labels. This thesis addresses the relational understanding of surgical activity by progressing from instrument perception to spatially grounded instrument--tissue interaction prediction.

\section{From Recognition to Grounded Interaction Understanding}
\label{sec:introduction-from-recognition}

Early work in surgical video analysis commonly separated perceptual and workflow tasks. Instrument detection and segmentation described the visible tools, while phase and action recognition described the progression or content of an operation. These tasks provide complementary information, but none alone gives a complete account of an instrument--tissue interaction. Tool recognition does not explain how a tool is being used; action recognition without spatial evidence does not identify the responsible object; and anatomical recognition without an associated action does not describe the event taking place.

Surgical action triplets provide a compact representation that connects these elements. An interaction is expressed as the \ivttriplet tuple, such as \texttt{<Grasper, Retract, Gallbladder>} or \texttt{<Hook, Dissect, Gallbladder>}. This formulation describes both an action and its anatomical context, and therefore captures finer-grained surgical meaning than instrument presence or workflow phase alone.

The CholecT series of datasets, particularly CholecT50, has been central to the development of surgical action triplet recognition \cite{NWOYE2022102433_rendevous}. TripNet and Rendezvous established influential modelling approaches for learning relationships between instruments, verbs, and targets \cite{1nywoye_rec_action_triplets,NWOYE2022102433_rendevous}. Their associated benchmarks enabled systematic comparison on fine-grained activity recognition in laparoscopic cholecystectomy. However, progress in recognising triplet labels did not by itself resolve where each interaction occurred.

Most surgical action triplet methods have historically treated the task as frame-level recognition. A model predicts that a triplet is present somewhere in an image, but does not necessarily associate it with a particular visible instrument. This distinction becomes important when a frame contains several tools, especially multiple instances of the same instrument class. If two graspers are visible, a frame-level annotation may indicate that a grasper retracts the liver and a grasper grasps the gallbladder, yet it does not explicitly identify which physical instance performs either interaction.

Spatial awareness therefore requires more than attaching a coarse location to a frame-level label. It requires an instance representation that distinguishes individual instruments and preserves the association between each instance and its predicted verb and target. Instance segmentation is well suited to this role because it separates overlapping or repeated tools while retaining their visible spatial extent. At the outset of this research, however, large-scale instance-level annotations for instruments in routine human laparoscopic surgery were scarce. This limited both the development of specialist instance segmentation models and the construction of strongly grounded interaction benchmarks.

This thesis connects these lines of research by progressing from instrument instance segmentation to grounded triplet prediction and, finally, interaction prediction for individually grounded instruments.

\section{Research Gap}
\label{sec:introduction-gap}

The literature review in \Chapref{chap:background} examines multitask learning as one important framework for surgical scene understanding. Its broader finding is that surgical understanding depends on multiple related perceptual, temporal, and semantic outputs. Multitask learning can exploit relationships between such outputs, but it is a means rather than the central subject of this thesis. The more fundamental question is whether the outputs needed for fine-grained interaction understanding are defined, annotated, and spatially connected in the first place.

The review identifies a progression from coarse workflow recognition towards increasingly fine-grained activity descriptions. However, surgical action triplets remained predominantly frame-level, while instrument instance segmentation had received substantially less attention than presence recognition, detection, or semantic segmentation. The field therefore lacked the combined supervision needed to connect semantically rich interactions to their responsible instrument instances.

These limitations define two connected gaps. First, a sufficiently large and diverse dataset of instrument instances is needed to support reliable instance segmentation in laparoscopic video. Second, action-triplet annotations must be aligned with those instances to establish a benchmark for grounded interaction understanding. Triplet labels without instance masks describe the interaction but leave its actor spatially ambiguous. Instance masks without verb and target labels localise the actor but do not describe its surgical role. Bringing these forms of supervision together is consequently a central contribution of this thesis.

Grounding also makes the prediction problem more demanding. A complete grounded triplet is correct only when the instrument instance is localised and its instrument, verb, and target components are all correctly identified. An error in any component can invalidate the complete prediction. Component-level results are therefore necessary for interpreting where performance is gained or lost, while complete-triplet performance measures whether the full relationship has been recovered. This conjunctive structure is particularly challenging for anatomical targets, which may be partially visible, visually similar, deformable, or defined through surgical context rather than local appearance alone.

The thesis addresses this target-prediction challenge through two complementary modelling directions. The first incorporates weak anatomical context into a strongly supervised instance segmentation architecture. The second separates instrument instance segmentation from semantic interaction prediction, allowing a multimodal large language model (MLLM) to predict \vtpair relationships for already segmented instruments. Together, these approaches investigate how spatial supervision, anatomical information, task factorisation, and temporal context can contribute to grounded surgical interaction understanding.

\section{Thesis Aim and Research Questions}
\label{sec:introduction-aim}

The aim of this thesis is to advance surgical scene understanding in minimally invasive surgery by developing datasets, benchmarks, and models for spatially grounded interpretation of instrument--tissue interactions in laparoscopic video.

This aim is addressed through the following research questions:

\begin{enumerate}
  \item How has surgical scene understanding progressed from individual perceptual tasks towards fine-grained interaction recognition, and what limitations prevent surgical action triplets from being spatially grounded to individual instrument instances?
  \item Can a large-scale instrument instance segmentation dataset be constructed for laparoscopic cholecystectomy by extending existing public surgical datasets?
  \item Can surgical action triplets be grounded to instrument instance masks, enabling spatially interpretable interaction predictions?
  \item Can weak anatomical priors improve target prediction and full triplet grounding when integrated into an instance segmentation architecture?
  \item Can an MLLM serve as a viable second-stage interaction predictor for grounded instrument instances, and how do alternative task formulations and sources of context affect its predictions?
\end{enumerate}

These questions form a cumulative programme of work, progressing from the identification of the research gap to instance-level spatial supervision, grounded triplet prediction, target-aware modelling, and two-stage multimodal interaction prediction.

\section{Contributions}
\label{sec:introduction-contributions}

The principal contributions of this thesis are summarised below.

\subsection{A Review-Based Framing of Surgical Scene Understanding}
\label{sec:introduction-contribution-review}

The thesis synthesises the literature on multitask learning in MIS and uses it to examine the wider development of surgical scene understanding. The review organises the field across perceptual, workflow, and activity-understanding tasks, and shows why multiple related outputs are needed to describe a surgical scene. It identifies the movement towards fine-grained action triplets, the limited spatial grounding available to triplet methods, and the relative scarcity of instance segmentation datasets. This analysis motivates the thesis's focus on connecting semantic interaction recognition with instance-level spatial evidence.

\subsection{CholecInstanceSeg}
\label{sec:introduction-contribution-cholecinstanceseg}

The thesis introduces CholecInstanceSeg, a large-scale open-access dataset for surgical instrument instance segmentation in laparoscopic cholecystectomy. It contains 41,933 annotated frames from 85 procedures and 64,483 instrument instances. The dataset provides semantic classes and instance identifiers, allowing individual instruments to be distinguished in scenes containing multiple or overlapping tools.

CholecInstanceSeg addresses the shortage of instance-level annotations in routine human laparoscopic surgery and supports the development and evaluation of specialist instrument segmentation models. More importantly for the thesis narrative, it provides the spatial substrate upon which later interaction labels can be grounded.

\subsection{CholecTriplet-Seg}
\label{sec:introduction-contribution-cholectripletseg}

Building on CholecInstanceSeg and CholecT50, the thesis introduces CholecTriplet-Seg, a dataset and benchmark for spatially grounded surgical action triplet understanding. It aligns instrument instance masks with verb and anatomical target labels across 30,955 frames and 49,866 grounded triplets from 50 laparoscopic cholecystectomy videos.

This contribution changes the triplet task from asking which interactions occur within a frame to asking which instrument instance performs each interaction. It thereby enables direct evaluation of the spatial association between an instrument and its surgical role, while retaining component-level and complete-triplet measures for diagnosing model behaviour.

\subsection{TargetFusionNet}
\label{sec:introduction-contribution-targetfusionnet}

The thesis proposes TargetFusionNet, a target-aware architecture for grounded surgical action triplet prediction. TargetFusionNet extends a strongly supervised instance segmentation framework with gated fusion of weak anatomical priors. The priors are treated as coarse contextual cues rather than ground-truth anatomical labels, allowing their contribution to target and triplet prediction to be assessed within the grounded benchmark.

Experiments on CholecTriplet-Seg demonstrate the importance of strong instance supervision and show that target-aware fusion improves complete-triplet prediction over the evaluated weakly grounded and strongly supervised baselines. They also establish that target prediction remains the principal bottleneck: improved instrument localisation and component predictions do not automatically yield proportionate gains in complete-triplet performance.

\subsection{Two-Stage Multimodal Interaction Prediction}
\label{sec:introduction-contribution-reasoning}

Finally, the thesis formulates surgical triplet segmentation as a two-stage problem. An instrument instance segmentation model first localises the visible instrument instances, after which an MLLM predicts the corresponding \vtpair interactions. Direct supervised fine-tuning establishes a viable second-stage baseline, while sequential dialogue prediction and short-range temporal context provide promising alternative formulations. Anatomical context, ontology guidance, longer temporal input, and distilled rationale supervision are also examined to determine whether additional structure or context consistently improves interaction prediction.

The complete pipeline is evaluated using predicted instrument instances, rather than only idealised ground-truth grounding. The results establish the viability of specialised grounding followed by multimodal interaction prediction and expose the effect of first-stage errors on the complete system. This chapter represents a forward-looking continuation of the dataset and grounded-modelling contributions, rather than replacing their central role in the thesis.

\section{Thesis Scope}
\label{sec:introduction-scope}

The thesis focuses on laparoscopic cholecystectomy, for which established public datasets make it possible to develop a coherent sequence of instance-segmentation, triplet-grounding, and interaction-prediction studies. The shared procedural setting allows the contributions to build upon Cholec80, CholecT50, and CholecSeg8k while maintaining consistent definitions of instruments, actions, and anatomical targets.

The thesis does not attempt to solve every aspect of surgical scene understanding. It does not provide exhaustive anatomical segmentation, model the complete temporal workflow of an operation, or validate a clinical decision-support system prospectively. Its focus is the more specific problem of connecting instrument instances to their actions and targets. Cholecystectomy provides the principal test bed; evaluating how well the formulation transfers to other procedures remains an important next step.

Within this scope, the thesis spans dataset construction, benchmark definition, specialist vision models, and multimodal modelling. These elements are unified by spatial grounding rather than by a commitment to any single learning paradigm. Multitask learning, anatomical priors, and multimodal language modelling are used where they help address particular parts of the grounded interaction problem.

\section{Thesis Structure}
\label{sec:introduction-structure}

The remainder of this thesis is organised as follows.

\Chapref{chap:background} reviews multitask learning in surgical scene understanding and examines the progression from individual perceptual tasks towards fine-grained activity recognition. It identifies the absence of instance-level spatial grounding as a key limitation of existing surgical action triplet formulations.

\Chapref{chap:cholecinstanceseg} presents CholecInstanceSeg, the spatial foundation of the thesis. It describes the data sources, annotation protocol, dataset construction, technical validation, and limitations of a large-scale instrument instance segmentation dataset for laparoscopic cholecystectomy.

\Chapref{chap:triplet-grounding} introduces CholecTriplet-Seg and TargetFusionNet. It defines triplet segmentation as an instance-grounded extension of surgical action triplet recognition, presents the aligned dataset and evaluation protocol, and assesses target-aware fusion for grounded prediction.

\Chapref{chap:two-stage-interaction} presents a two-stage formulation comprising instrument instance segmentation followed by MLLM-based interaction prediction. It establishes direct supervised fine-tuning, evaluates several task formulations and contextual adaptations, and assesses the complete pipeline using predicted instrument instances.

\Chapref{chap:discussion-conclusion} synthesises the thesis findings across the review, datasets, benchmarks, and models. It discusses the difficulty and evaluation of grounded triplet prediction, identifies the principal limitations of the work, and outlines future research towards robust temporal modelling and evaluation across surgical procedures.
